% ============================================================
% OP3_1a_bridge_recon_v3.tex  (PROPOSAL-ONLY; the V-coupling
% computation for M-(alpha,beta), per GPT's go-ahead and warning
% that the potential sector may be a constraint/failure mode
% rather than a bridge). NO preprint edits, NO mechanism claims.
% All statements are for the 1D (profile) reduction; 4D readings
% are flagged as interpretation.
% ============================================================
\section*{OP3.1a bridge, v3: the potential-sector route computed
(proposal-only)}

\subsection*{0. Setup}
$S=\int d^4X\,[\,P(X)+V(\Phi)\,]$, $X=\partial_A\Phi\,
\partial_A\Phi$, with the recorded P3 potential
$V=-\tfrac{\mu^2}{2}\Phi^2+\tfrac{\lambda}{4}\Phi^4$
(well bottom $\phi_0$, $V''(\phi_0)=m_\sigma^2=2\mu^2$).
Profile reduction: $\Phi=\Phi(s)$ along a fixed direction,
$X=\Phi'^2$.
The Euler--Lagrange equation is
$2\,\tfrac{d}{ds}\bigl(P_X\,\Phi'\bigr)=V'(\Phi)$.

\subsection*{1. Lemma B1 (first integral and a structural
identity)}
Define $Q(X):=2X\,P_X(X)-P(X)$.
The profile equation has the first integral
\[
Q(X)\;=\;E+V(\Phi),\qquad E=\mathrm{const},
\]
(the 1D ``energy'' $2P_X\Phi'^2-P-V$).
Moreover
\[
Q'(X)=P_X(X)+2X\,P_{XX}(X),
\]
i.e.\ \emph{$Q'$ is exactly the longitudinal fluctuation
coefficient} of the computed mini-lemma.
Along an orbit, $dX/d\Phi=V'(\Phi)/Q'(X)$: the gradient magnitude
is slaved to the potential wherever $Q'\neq0$.

\subsection*{2. Lemma B2 (with the recorded $V$, exact homogeneous
gradient order is excluded; deformation and patch scale)}
For the linear profile, $\tfrac{d}{ds}(P_X\Phi')=0$ while
$V'(\bar\Phi(s))\neq0$ away from $\Phi\in\{0,\pm\phi_0\}$;
hence the exact constant-gradient background of the pure-$P$
motivation does not survive the recorded $V$.
Near a well bottom ($\Phi=\phi_0+u_*s$, $V'\approx m_\sigma^2
u_*s$), the first integral gives the deformation law
\[
X(s)-u_*^2\;\approx\;\frac{\mu^2u_*^2\,s^2}{Q'(u_*^2)},
\qquad Q'(u_*^2)=2u_*^2P_{XX}(u_*^2)>0
\ \text{at the recorded point},
\]
so the gradient magnitude departs from $u_*^2$ by order one over
the patch length
\[
\ell_V\;\sim\;\frac{\sqrt{Q'(u_*^2)}}{\mu}\;\sim\;m_\sigma^{-1}
\times O\!\bigl(\sqrt{Q'}\bigr).
\]
\emph{Hedged observation (1D; interpretive in 4D):} the
potential-sector coupling limits quasi-homogeneous gradient order
to patches of the order of the recorded condensate correlation
scale $\xi_{\rm cond}\sim m_\sigma^{-1}$ --- resonating with, but
not derived to be, the recorded coarse-graining regime
(eq:cg\_regime).

\subsection*{3. Lemma B3 (the potential sector contributes no
derivative terms: not a kinetic bridge)}
Around any background, the quadratic action is
\[
S^{(2)}=\int d^4X\Bigl[\,P_X(\bar X)\,(\partial\chi)^2
+2P_{XX}(\bar X)\,(\partial^A\bar\Phi\,\partial_A\chi)^2
+\tfrac12 V''(\bar\Phi)\,\chi^2\Bigr],
\]
i.e.\ $V$ enters only through the position-dependent mass
$V''(\bar\Phi)=-\mu^2+3\lambda\bar\Phi^2$ (positive
$m_\sigma^2$ near the wells, a tachyonic band
$|\bar\Phi|<\phi_0/\sqrt3$ between them) and through the
background deformation of Lemma~B2; the kinetic tensor structure
is unchanged.
The same holds for mixed $P(X,\Phi)$ at this order: the cross
term $P_{X\Phi}\,\delta X\,\chi\supset 2u\,P_{X\Phi}\,
(n\!\cdot\!\partial\chi)\chi=u\,P_{X\Phi}\,n\!\cdot\!\partial
(\chi^2)$ reduces by parts to a mass-sector contribution; the
kinetic coefficients remain $P_X$ and $2XP_{XX}$ evaluated at the
local background point.
\emph{Verdict:} the bridge candidate ``coupling to the potential
sector'' cannot resolve the kinetic part of M-($\alpha,\beta$);
its computed r\^ole is (i)~a background constraint (Lemma~B2) and
(ii)~mass-sector structure --- a constraint/failure mode for the
global linear route, as anticipated in review.

\subsection*{4. Observation B4 (sign chain; the sharpened form of
M-($\alpha,\beta$))}
Matching the local kinetic tensor to
$\beta\delta^{AB}-\alpha n^An^B$ gives
$\beta=P_X(\bar X)$ and
$\beta-\alpha=P_X(\bar X)+2\bar X P_{XX}(\bar X)=Q'(\bar X)$.
Hence
\[
\alpha>\beta>0
\quad\Longleftrightarrow\quad
P_X(\bar X)>0\ \text{and}\ Q'(\bar X)<0 .
\]
But $Q'<0$ is precisely the regime where the first-integral
slaving $dX/d\Phi=V'/Q'$ sits on a descending (1D-unstable)
branch of $Q$: the potential-sector mechanics cannot itself
stabilise the operative point there.
\emph{Reading (hedged):} the Lorentzian-phase tensor requires the
operative background to be, in the 1D sense, on an unstable
longitudinal branch --- consistent with the recorded
saddle-driven-restructuring character, and showing that the
bridge must invoke structure beyond 1D $P{+}V$ mechanics
(orientation/EFT-level stiffness, constraint sector, or the level
distinction between bare and EFT tensors).
No inconsistency is asserted; this sharpens where the bridge can
and cannot live.

\subsection*{5. Updated bridge-candidate ledger}
Potential-sector coupling: \emph{reclassified} --- constraint
C-V (Lemma~B2) + mass sector (Lemma~B3); not a kinetic bridge.
Mixed $P(X,\Phi)$: kinetic-tensor signs still governed by the
local $(P_X,P_{XX})$; can relocate the operative point
($P_X\neq0$ via the first integral) but inherits the B4 sign
chain.
Live candidates for the kinetic part: orientation/EFT-level
stiffness ($\mathcal C_n$ as the recorded transverse supplier);
constraint-sector or composite-variable effects; explicit
bare-vs-EFT level distinction.
Selection/stabilisation at $Q'<0$: belongs with the recorded
saddle framework and OP3.1c, not with 1D mechanics.

\subsection*{6. What is NOT claimed}
No 4D theorems (all computations are profile-reduction
statements); no claim that the recorded $P(X)$ motivation is
wrong; no claim that OP-grad has failed --- only that one bridge
candidate is reclassified as a constraint; no insertion; no
status changes; no numbers for $u_*(\mu,\lambda)$.

\subsection*{7. Questions for GPT}
(Q1)~Approve Lemmas B1--B3 and Observation B4 as computed (1D)
statements?
(Q2)~Eventual insertion shape (after your sign-off): a compact
``Potential-sector constraint (computed)'' paragraph next to the
implemented OP3.1a block --- first integral $Q(X)=E+V$, the
identity $Q'=$ longitudinal coefficient, the exclusion of the
exact linear profile under the recorded $V$ with the patch scale,
the no-derivative-terms verdict, and the B4 equivalence
$\alpha>\beta>0\Leftrightarrow P_X>0,\ Q'<0$ as the sharpened
matching question?
(Q3)~The $\xi_{\rm cond}$ resonance of B2: include as a hedged
one-sentence observation, or proposal-only?
(Q4)~Next computation: the orientation-sector route (does the
recorded $\mathcal C_nu^2(\partial n)^2$ supply $\beta>0$ with the
required sign structure at EFT level?), or pause the OP3.1a column
here?
